\clearpage
\setcounter{page}{2}
\pagestyle{plain}
\singlespacing

{\crsTOCheading СОДЕРЖАНИЕ\par\vspace{1ex}}

\crsTOCline[0pt]{\normalfont\textbf{ВВЕДЕНИЕ}}{3}
\crsTOCline[0pt]{\normalfont\textbf{ГЛАВА\,1}. Теоретические основы генерации случайных чисел}{4}
\crsTOCline[1.5em]{\normalfont{}1.1. Основные понятия и определения}{4}
\crsTOCline[1.5em]{\normalfont{}1.2. Классификация генераторов случайных чисел}{5}
\crsTOCline[1.5em]{\normalfont{}1.3. Методы оценки случайности последовательностей}{7}
\crsTOCline[0pt]{\normalfont\textbf{ГЛАВА\,2}. Анализ физических источников энтропии и их применение в генераторах истинно случайных последовательностей}{11}
\crsTOCline[1.5em]{\normalfont{}2.1. Электронные источники шума}{11}
\crsTOCline[1.5em]{\normalfont{}2.2. Оптические и квантовые источники}{17}
\crsTOCline[1.5em]{\normalfont{}2.3. Механические и акустические источники}{20}
\crsTOCline[1.5em]{\normalfont{}2.4. Постобработка}{23}
\crsTOCline[1.5em]{\normalfont{}2.5. Уязвимости и атаки}{25}
\crsTOCline[0pt]{\normalfont\textbf{ГЛАВА\,3}. Практический анализ и сравнение источников энтропии}{29}
\crsTOCline[1.5em]{\normalfont{}3.1. Сравнительный анализ источников энтропии}{29}
\crsTOCline[1.5em]{\normalfont{}3.2. Архитектура экспериментального стенда и программная инфраструктура}{30}
\crsTOCline[1.5em]{\normalfont{}3.3. Натурная оценка теплового шума резистора по схеме с двумя неинвертирующими каскадами (три эксперимента)}{33}
\crsTOCline[1.5em]{\normalfont{}3.4. Экспериментальная оценка акустического источника энтропии на базе электретного микрофона}{38}
\crsTOCline[1.5em]{\normalfont{}3.5. Перспективные направления развития физических генераторов случайных последовательностей (ФГСП)}{46}
\crsTOCline[0pt]{\normalfont\textbf{ЗАКЛЮЧЕНИЕ}}{48}
\crsTOCline[0pt]{\normalfont\textbf{СПИСОК ИСПОЛЬЗОВАННЫХ ИСТОЧНИКОВ}}{50}
\crsTOCline[0pt]{\normalfont\textbf{ПРИЛОЖЕНИЯ}}{52}
\crsTOCline[1.5em]{\normalfont Приложение~1\,--- Статистические тесты пакета NIST STS}{52}
\crsTOCline[1.5em]{\normalfont Приложение~2\,--- Сравнительная таблица характеристик источников энтропии}{55}
